\chapter{Introduction}

This is the beginning. I bet you wish that you had very started before long. I
shall now write two paragraphs. This is the second paragraph. In recent times
the way in which the human race uses and generates its energy has become of
extreme social, economical and environmental importance. Predictions of damaging
increases in mean global temperatures \citep[see][]{Solomon:2007:CUP} has
increased the need for carbon emitting energy technologies to be replaced by low
carbon alternatives. Additional long term economic factors are also playing
their part in shifting momentum to new technologies, in particular the concerns
about oil and natural gas supplies. Demand for oil and gas is expected to
outstrip supply within the current century, leading to inflated prices and
energy security issues. Unfortunately, although the resource and environmental
issues are occurring simultaneously, they are not necessarily mutually
supportive. For instance, once the price for oil has reached a certain value it
can be economically synthesised from coal, providing no environmental benefits.
In fact, if the market were left to choose a method for replacing dwindling oil
resource then this is one of the most likely substitutes; even in light of
population growth, coal reserves are estimated to last for hundreds of years
\citep[see][]{Jaccard:2005:CUP} and thus the cost is low.

With markets failing to deliver the necessary changes, it has become the
responsibility of governments to intervene under the premise that the predicted
environment and economical consequences of inaction will outweigh the costs
incurred by immediate action. This is a tough political task as, particularly in
the UK, energy markets have followed a trend of liberalised trading which makes
strategic decision making extremely challenging. Private investment is unlikely
to match policy, unless the policy is seen to be long term and economically
advantageous; a virtual impossibility when scientific and economic opinions of
the market requirements are so uncertain. In addition, the likelihood of direct
governmental intervention in the form of capital investment is now looking more
unlikely, due to the recent financial crisis. Capital spending is set to fall
significantly in order to reduce deficits endured to rescue the banking sector.

Despite the potential funding problems, the United Kingdom (UK) Government's
Climate Change Act 2008 \citep*[see][2008]{CCA:2008:Defra} introduced for the
first time legally binding targets for greenhouse gas emissions from within the
UK. The targets set an 80\% reduction of greenhouses gases by 2050 and a 26\%
reduction of carbon-dioxide $\text{CO}_{2}$ emissions by 2020 with respect to
1990 levels. This is set in the context of Britain's pre-existing target to
reduce emissions by 12\% as a ratified signatory of the Kyoto Protocol
and a European Union proposal to cut EU wide greenhouse gas emissions by 20\% by
2020 each with respect to 1990 levels. Electricity generation in the United
Kingdom accounts for 37\% of all emissions and as part of the
government's greenhouse gas reduction strategy it will seek to reduce this to
zero by 2030 \citep[see][2008]{CCC:2008}. Given the enormity of this task it is
foreseen that a varied mix of low/zero carbon electricity generating
technologies will be required.
\begin{quote}
 ``It is a well known fact that the use of a high-order panel method is more
accurate than the
low-order panel method or the discrete vortex method in computing the velocity
field as the
appearance of instabilities in the vortex sheet due to the spurious numerical
effects introduced by
a too-crude representation.''
\end{quote}
As such, a melting pot of renewable and sustainable energy technologies have
begun to compete to become part of a `post-carbon' energy mix. Some of these
technologies are well established, such a nuclear, biomass and wind, others less
so, such as solar photovoltaic and marine energy. In addition, carbon emitting
technologies such as coal remain an attractive option for governments and
investors alike, as carbon capture and storage technology promises to store away
the greenhouse gases emitted in combustion. However, carbon capture and storage
is still an unproven technology and thus, the future costs of this and many
other of the new energy technologies are hard to predict and are, at present,
highly contestable; ultimately, over the long term, energy cost is the most
likely deciding factor as to which technologies will feature most. Hence, each
industry is working hard to reduce costs via research and innovation.

\begin{figure}[htb]
\centering
\includegraphics[width=0.3\textwidth]{bottomkutta.eps}
\caption{Alternative orientations of the wake vortex sheet as it separates from
the trailing edge of the aerofoil.}
\label{fig:bottomkutta}
\end{figure}